\documentclass[10pt,a4paper]{amsart}
\usepackage[margin=19mm]{geometry}
\usepackage{amsmath,amssymb,mathtools}
\usepackage[scr=rsfs,cal=euler]{mathalfa}
\usepackage{xfrac}
\usepackage[unicode,hidelinks,bookmarks=false]{hyperref}
\newcommand{\ie}{i.e.\ }
\newcommand{\cf}{cf.\ }
\newcommand{\resp}{respectively\ }
\newcommand{\Subsection}{\subsection}
\providecommand{\nobreakdash}{\nobreak\hbox{-}}
\setlength{\emergencystretch}{2em}
\multlinegap0pt
\usepackage{xcolor}
\usepackage{csquotes}
\usepackage{mathtools}
\usepackage{bm}
\usepackage{amssymb}
\newtheorem{lemma}{Lemma}
\usepackage{cleveref}
\crefname{lemma}{Lemma}{Lemmas}
\title{Rerouting Curves on Surfaces}
\author{Timo Brand, Stefan Felsner, Henry F\"orster, Stephen Kobourov, Anna Lubiw, Yoshio Okamoto, J\'anos Pach, Csaba D. T\'oth, G\'eza T\'oth, Torsten Ueckerdt,, Pavel Valtr}
\date{Source: arXiv:2607.05362v1 (CC BY 4.0)}
\begin{document}
\maketitle

\noindent\emph{Source and licence.} Rerouting Curves on Surfaces. Version arXiv:2607.05362v1 (CC BY 4.0), distributed by arXiv under the Creative Commons Attribution 4.0 International licence, http://creativecommons.org/licenses/by/4.0/, as recorded in the arXiv licence metadata for this exact version and shown on the abstract page https://arxiv.org/abs/2607.05362v1. Full licence notice: \texttt{licenses/arxiv-2607.05362-CC-BY-4.0.txt}.\par\smallskip

\noindent\textit{Editorial scope.} Counted unit: the article's lemma lem:hug-a-tree (source0.tex lines 1075--1080). The article's complete original proof of it is delivered with it (source0.tex lines 1081--1113, one \texttt{proof} environment of 33 lines). No mathematics is written, repaired or re-derived: the statement and the proof are the article's own lines, in the article's own notation, and no retained line is substituted or re-flowed.  No further source-local context is needed: every cross-reference of every retained line resolves inside the two delivered spans.  One declared adaptation: exactly 1 ancillary strings inside the retained spans are omitted (image-inclusion commands and, where the source repeats a label, the repeated label definitions) and no other line differs from the source; the figure environments, with their placement, captions and labels, are kept, and the package records the omitted strings in fidelity receipts bound to the affected bodies.  No retained line contains a citation, so the wrapper carries no bibliography and no bibliography entry.  The retained lines contain the article's own diagram code, which the wrapper typesets with the standard tikz packages; no image file is delivered and the package declares no asset dependency.  Editorial formatting only, in addition: an AMS \texttt{amsart} preamble that restates the article's own declarations and macros used by the retained lines, namely the environments \texttt{lemma} on separate counters, together with the standard packages those lines need.  The retained environments print in this document as Lemma 1 in order of appearance, whereas the article prints the same items with the numbering of its own full document.  The complete native source of the article as distributed in the arXiv e-print for this exact version is bundled unchanged as \texttt{sources/204538/source0.tex}.
\begin{lemma}\label{lem:hug-a-tree}
    Let $\mathcal{B}$ be a $\Sigma$-embedding of a connected planar graph $G = (V,E)$ on an orientable surface $\Sigma$ of genus $g\geq 1$ so that $\mathcal{B}$ has the same rotation system as a plane embedding $\mathcal{E}$ of $G$.
    Let $f$ be a face of the plane embedding $\mathcal{E}$, and let $T$ be a spanning tree of $G$.
    Then $\mathcal{B}$ can be reconfigured to an embedding $\mathcal{B}'$ in
    an $\varepsilon$-strip system of $\mathcal{B}(T)$ so that the outer face of $\mathcal{B}'$ corresponds to $f$, and the rotation system remains the same throughout the reconfiguration sequence.
\end{lemma}

\begin{proof}
    Let $\varepsilon>0$ be sufficiently small so that
    $\mathcal{B}(T)$ admits an $\varepsilon$-strip system, $U_\varepsilon$ intersects an edge segment of $\mathcal{B}$ if and only if it is part of $\mathcal{B}(T)$ or incident to a vertex in $V$; and its intersection with any such edge segment is connected.
    Recall that $U_\varepsilon$ is a neighborhood of $\mathcal{B}(T)$
    which contains the disks and strips of the $\varepsilon$-strip system.
    In particular, $U_\varepsilon$ is homeomorphic to a disk.
    Similarly, we can define a $\delta$-strip system for $\mathcal{E}(T)$ in the plane, and let $U_\delta$ be a neighborhood of $\mathcal{E}(T)$.
    Since $\mathcal{B}$ and $\mathcal{E}$ have the same rotation system, there is a bijection $\varphi:\partial U_\varepsilon\to \partial U_\delta$
    that maps the crossings $\mathcal{B}(e)\cap \partial U_\varepsilon$ to the crossings $\mathcal{E}(e)\cap \partial U_\delta$ for all $e\in E$.

    We may assume, without loss of generality, that $f$ is the outer face in $\mathcal{E}$. Every edge $e\in E$ that is not in $T$ determines a unique cycle $C_e$ in $T\cup\{e\}$, and its embedding $\mathcal{E}(C_e)$ is a Jordan curve that encloses one or more bounded faces in $\mathcal{E}$. For every $e\in E\setminus E(T)$, let ${\rm rank}(e)$ be the number of bounded faces of $\mathcal{E}$ enclosed by $\mathcal{E}(C_e)$. (Note, however, that $\mathcal{B}(C_e)$ is not necessarily a separating cycle in $\Sigma$.)


    \begin{figure}[htbp]
        \centering
        
        \caption{Rerouting an edge $e\in E\setminus E(T)$ to an edge curve in the $\varepsilon$-strip system of $T$.}
        \label{fig:hug-a-tree}
    \end{figure}

    We describe an algorithm that reconfigures $\mathcal{B}$ to an embedding $\mathcal{B}'$ contained in $U_\varepsilon$ such that its outer face w.r.t.\ $U_\varepsilon$ corresponds to face $f$ (of the plane embedding $\mathcal{E}$). We reconfigure every edge $e\in E\setminus E(T)$ in increasing order by rank (ties are broken arbitrarily).

    Assume that we have already rerouted some of the edges, and $e\in E\setminus E(T)$ is the next edge to handle; see \Cref{fig:hug-a-tree}.
    We may assume, without loss of generality, that $\mathcal{B}$ is the current embedding, and $\mathcal{B}$ embeds all edges of rank less than ${\rm rank}(e)$ in the interior of $U_\varepsilon$.
    For every edge $e'\in E\setminus E(T)$, if $\mathcal{E}(e')$
    lies in the interior of $\mathcal{E}(C_e)$, then ${\rm rank}(e')<{\rm rank}(e)$, and by assumption $\mathcal{B}(e')$ is already in the interior of $U_\varepsilon$.
    By the choice of $\varepsilon>0$, $\mathcal{B}(e)$ intersects $\partial U_\varepsilon$ in two points, which partition the closed curve $\partial U_\varepsilon$ into two arcs. Let $\gamma$ be the arc of $\partial U_\varepsilon$ for which $\varphi(\gamma)$ is in the interior of $\mathcal{E}(C_e)$. Note that $\gamma$ cannot cross any edge curve of $\mathcal{B}$, since $\varphi$ would map any crossing $\mathcal{B}(e')\cap \gamma$ to the interior of $\mathcal{E}(C_e)$, which would imply that ${\rm rank}(e')<{\rm rank}(e)$, and $\mathcal{B}(e')$ lies in the interior of $U_\varepsilon$.

    Consequently, we can now reroute $e$ in the $\varepsilon$-strip system closely following the arc $\gamma$ between the two intersection points in $\mathcal{B}(e)\cap \partial U_\varepsilon$. Clearly, by rerouting edge $e$, we did not change the rotation system.

    After successively rerouting all edges in $E\setminus E(T)$, let $\mathcal{B}'$ denote the resulting $\Sigma$-embedding of $G$.
    Note that all bounded faces of $\mathcal{B}'$ lie in the interior of $U_\varepsilon$.
\end{proof}

\end{document}
